\documentclass[11pt]{article}
% ======================================================================
%  examstyle.tex — EPFL-style exam layout for the weekly Time Series exams
%  Shared by every Exam_Week_NN(.tex / _Solutions.tex). Compiles with pdflatex.
% ======================================================================
\usepackage[utf8]{inputenc}
\usepackage[T1]{fontenc}
\usepackage[english]{babel}
\usepackage[a4paper,margin=2.4cm]{geometry}
\usepackage{amsmath,amssymb,amsthm,mathtools}
\usepackage{bm}
\usepackage{enumitem}
\usepackage{array}

\setlength{\parindent}{0pt}
\setlength{\parskip}{0.5em}
\emergencystretch=3em
\allowdisplaybreaks[1]

% ---------- math macros (same names as the rest of the project) ----------
\newcommand{\E}{\mathbb{E}}
\DeclareMathOperator{\Var}{Var}
\DeclareMathOperator{\Cov}{Cov}
\DeclareMathOperator{\Corr}{Corr}
\DeclareMathOperator*{\argmin}{arg\,min}
\DeclareMathOperator{\MSE}{MSE}
\DeclareMathOperator{\trace}{tr}
\newcommand{\Reals}{\mathbb{R}}
\newcommand{\Ints}{\mathbb{Z}}
\newcommand{\Comp}{\mathbb{C}}
\newcommand{\Nat}{\mathbb{N}}
\newcommand{\Prob}{\mathbb{P}}
\newcommand{\Tr}{^{\mathsf{T}}}
\newcommand{\conj}[1]{\overline{#1}}
\newcommand{\abs}[1]{\left\lvert #1 \right\rvert}
\newcommand{\norm}[1]{\left\lVert #1 \right\rVert}
\newcommand{\ind}{\mathbf{1}}
\newcommand{\eps}{\varepsilon}
\newcommand{\diff}{\nabla}
\newcommand{\acvs}{\gamma}
\newcommand{\acf}{\rho}
\newcommand{\pacf}{\alpha}
\newcommand{\sdf}{\mathcal{S}}
\newcommand{\Bsh}{B}
\newcommand{\iid}{\overset{\text{iid}}{\sim}}

\setlist[enumerate]{leftmargin=2em,topsep=2pt,itemsep=4pt}
\setlist[itemize]{leftmargin=1.6em,topsep=2pt,itemsep=2pt}

\newcounter{exo}

% ---------- exam cover header :  \examheader{exam title}{duration} ----------
\newcommand{\examheader}[2]{%
  \thispagestyle{empty}
  \begin{center}
    {\Large\bfseries Time Series}\\[3pt]
    Spring Semester 2026, EPFL\\
    \'Ecole Polytechnique F\'ed\'erale de Lausanne\\
    Prof.\ Dr.\ Sofia C.\ Olhede\\[7pt]
    \hrule height 0.8pt
    \vspace{9pt}
    {\large\bfseries Practice Exam --- #1}\\[4pt]
    {\normalsize Duration: #2 \quad$\cdot$\quad No calculator \quad$\cdot$\quad Answer on paper}\\
    \vspace{8pt}
    \hrule height 0.8pt
  \end{center}
  \vspace{14pt}
  \begin{tabular}{@{}p{8.2cm}@{}p{8cm}@{}}
    Name:\enspace\hrulefill & Surname:\enspace\hrulefill
  \end{tabular}\\[14pt]
  SCIPER (Student Card Number):\enspace\hrulefill
  \vspace{16pt}
}

% ---------- points table (fixed: 4 problems) ----------
\newcommand{\pointstable}{%
  \begin{center}
  \renewcommand{\arraystretch}{1.5}
  \begin{tabular}{|p{5cm}|p{4cm}|}
    \hline
    \textbf{Exercise} & \textbf{Points}\\\hline
    1 & \\\hline
    2 & \\\hline
    3 & \\\hline
    4 & \\\hline
    \textbf{Total} & \\\hline
  \end{tabular}
  \end{center}
}

% ---------- problem heading (exam: one problem per page) ----------
\newcommand{\problem}[1]{%
  \clearpage
  \stepcounter{exo}%
  \begin{center}\textbf{\large Problem \theexo\ (#1 points)}\end{center}
  \vspace{4pt}
}
% ---------- answer space: fill the statement page + one blank answer page ----------
% Put \answerpage at the END of each problem so every exercise gets its own page(s).
\newcommand{\answerpage}{\par\vfill\newpage\null\newpage}

% ---------- solutions document ----------
\newcommand{\solheader}[1]{%
  \begin{center}
    {\Large\bfseries Time Series --- Solutions}\\[3pt]
    {\large Practice Exam --- #1}\\[2pt]
    EPFL, MATH-342 \quad$\cdot$\quad Prof.\ S.\ C.\ Olhede\\[5pt]
    \hrule height 0.8pt
  \end{center}
  \vspace{10pt}
}
\newcommand{\solproblem}[1]{%
  \bigskip
  \stepcounter{exo}%
  {\large\bfseries Problem \theexo\ (#1 points)}\par\nobreak\vspace{2pt}
}
% Rappel de l'énoncé avant la solution (dans les corrigés)
\newcommand{\statementstart}{\par\vspace{1pt}{\bfseries\itshape Statement.}\par\nobreak\begingroup\small\itshape}
\newcommand{\statementend}{\par\endgroup\vspace{5pt}\hrule\vspace{6pt}{\bfseries Solution.}\par\nobreak\vspace{2pt}}

\begin{document}

\solheader{Week 4: Parametric Estimation --- Yule--Walker \& Least Squares}

% ---------------------------------------------------------------
\solproblem{5}
\statementstart
Let $\{X_t\}$ be a mean-zero causal $\mathrm{AR}(2)$ process
$X_t=\phi_1X_{t-1}+\phi_2X_{t-2}+\eps_t$.

\textbf{(a)} Starting from the model equation, derive the \emph{Yule--Walker equations}: multiply by $X_{t-k}$, take expectations, and explain precisely where \emph{causality} is used to kill the noise cross-term. State the resulting linear system $\boldsymbol\acvs=\Gamma\boldsymbol\phi$ and the variance identity for $\sigma^2$. \hfill(2 pts)

\textbf{(b)} A series gives the sample autocovariances $\hat\acvs_0=5$, $\hat\acvs_1=3$, $\hat\acvs_2=1$. Compute the Yule--Walker estimates $\hat\phi_1,\hat\phi_2$ and $\hat\sigma^2$. \hfill(2 pts)

\textbf{(c)} Verify that the fitted model of part (b) is causal (stationary) by locating the roots of its AR polynomial. \hfill(1 pt)
\statementend

\textbf{(a)} \emph{Derivation.} Multiply the model
$X_t=\phi_1X_{t-1}+\phi_2X_{t-2}+\eps_t$ by $X_{t-k}$ and take expectations. With
mean zero, $\E[X_{t-j}X_{t-k}]=\acvs_{k-j}$, so for every $k\ge0$,
\[
\acvs_k=\phi_1\acvs_{k-1}+\phi_2\acvs_{k-2}+\E[\eps_tX_{t-k}].
\]
Because $\{X_t\}$ is \emph{causal}, $X_{t-k}$ is a one-sided combination of
$\eps_{t-k},\eps_{t-k-1},\dots$ (innovations \emph{strictly earlier} than
$\eps_t$ when $k\ge1$), and white noise is uncorrelated across time, so
$\E[\eps_tX_{t-k}]=0$ for $k\ge1$. This is the single place where causality
enters. Taking $k=1,2$ gives, with $\acvs_{-1}=\acvs_1$,
\[
\boldsymbol\acvs=\Gamma\boldsymbol\phi,\qquad
\begin{pmatrix}\acvs_1\\ \acvs_2\end{pmatrix}
=\begin{pmatrix}\acvs_0 & \acvs_1\\ \acvs_1 & \acvs_0\end{pmatrix}
\begin{pmatrix}\phi_1\\ \phi_2\end{pmatrix}.
\]
For $k=0$ the cross-term does \emph{not} vanish: multiplying instead by $X_t$,
$\E[\eps_tX_t]=\E[\eps_t(\phi_1X_{t-1}+\phi_2X_{t-2}+\eps_t)]=\sigma^2$
(all $\E[\eps_tX_{t-j}]=0$ for $j\ge1$), giving the variance identity
\[
\acvs_0=\phi_1\acvs_1+\phi_2\acvs_2+\sigma^2,
\qquad\text{i.e.}\qquad
\sigma^2=\acvs_0-\phi_1\acvs_1-\phi_2\acvs_2 .
\]

\textbf{(b)} With $\hat\acvs_0=5,\ \hat\acvs_1=3,\ \hat\acvs_2=1$ the
Yule--Walker system is
\[
\begin{pmatrix}5 & 3\\ 3 & 5\end{pmatrix}
\begin{pmatrix}\hat\phi_1\\ \hat\phi_2\end{pmatrix}
=\begin{pmatrix}3\\ 1\end{pmatrix}.
\]
The inverse of $\left(\begin{smallmatrix}a&b\\ b&a\end{smallmatrix}\right)$ is
$\tfrac{1}{a^2-b^2}\left(\begin{smallmatrix}a&-b\\ -b&a\end{smallmatrix}\right)$;
here $a=5,\ b=3$, determinant $a^2-b^2=25-9=16$, so
\[
\hat{\boldsymbol\phi}
=\frac{1}{16}\begin{pmatrix}5 & -3\\ -3 & 5\end{pmatrix}\begin{pmatrix}3\\ 1\end{pmatrix}
=\frac{1}{16}\begin{pmatrix}15-3\\ -9+5\end{pmatrix}
=\frac{1}{16}\begin{pmatrix}12\\ -4\end{pmatrix}
=\begin{pmatrix}\tfrac34\\[2pt] -\tfrac14\end{pmatrix}.
\]
So $\hat\phi_1=0.75,\ \hat\phi_2=-0.25$. The variance estimate is
\[
\hat\sigma^2=\hat\acvs_0-\hat\phi_1\hat\acvs_1-\hat\phi_2\hat\acvs_2
=5-\tfrac34(3)-\big(-\tfrac14\big)(1)
=5-\tfrac94+\tfrac14=5-2=3 .
\]
Hence $\hat\phi_1=\tfrac34,\ \hat\phi_2=-\tfrac14,\ \hat\sigma^2=3$.

\textbf{(c)} The fitted AR polynomial is
$\Phi(z)=1-\tfrac34 z+\tfrac14 z^2$. Solving $\tfrac14 z^2-\tfrac34 z+1=0$,
i.e.\ $z^2-3z+4=0$,
\[
z=\frac{3\pm\sqrt{9-16}}{2}=\frac{3\pm\sqrt{-7}}{2}=\tfrac32\pm\tfrac{\sqrt7}{2}\,i,
\qquad
\abs{z}=\sqrt{\Big(\tfrac32\Big)^2+\Big(\tfrac{\sqrt7}{2}\Big)^2}=\sqrt{\tfrac94+\tfrac74}=\sqrt4=2 .
\]
Both roots have modulus $2>1$, i.e.\ lie outside the unit circle, so the fit is
\textbf{causal (stationary)}. (Equivalently, the AR(2) triangle:
$\phi_1+\phi_2=\tfrac12<1$, $\phi_2-\phi_1=-1<1$, $\abs{\phi_2}=\tfrac14<1$.)

% ---------------------------------------------------------------
\solproblem{5}
\statementstart
A series of length $n=150$ is modelled as a mean-zero $\mathrm{AR}(3)$ process
$X_t=\phi_1X_{t-1}+\phi_2X_{t-2}+\phi_3X_{t-3}+\eps_t$. The estimated
autocorrelations are
\[
\hat\acf_1=0.8,\qquad \hat\acf_2=0.6,\qquad \hat\acf_3=0.4,
\]
and the sample variance is $\hat\acvs_0=10$.

\textbf{(a)} Explain why the Yule--Walker fit can be written entirely with the \emph{correlation} matrix, $\hat{\boldsymbol\phi}=\bar\Gamma_3^{-1}\hat{\boldsymbol\acf}_3$, i.e.\ why $\hat\acvs_0$ cancels in the system for $\boldsymbol\phi$. Write out the $3\times3$ system explicitly. \hfill(1.5 pts)

\textbf{(b)} Solve the system for $\hat\phi_1,\hat\phi_2,\hat\phi_3$ (e.g.\ by Gaussian elimination) and verify your solution by substituting it back into the three equations. \hfill(2.5 pts)

\textbf{(c)} Compute the innovation-variance estimate $\hat\sigma^2$ using $\hat\acvs_0=10$. \hfill(1 pt)
\statementend

\textbf{(a)} The population Yule--Walker system is
$\boldsymbol\acvs=\Gamma\boldsymbol\phi$, with $\Gamma=(\acvs_{i-j})$ and
$\boldsymbol\acvs=(\acvs_1,\acvs_2,\acvs_3)\Tr$. Factor out $\acvs_0$: writing
$V=\acvs_0^{-1}I$ and multiplying both sides by $V$,
\[
V\boldsymbol\acvs=V\Gamma\boldsymbol\phi
\;\Longrightarrow\;
\hat{\boldsymbol\acf}_3=\bar\Gamma_3\,\boldsymbol\phi,
\]
where every entry has been divided by $\acvs_0$, so
$\bar\Gamma_3=(\hat\acf_{i-j})$ and $\hat{\boldsymbol\acf}_3=(\hat\acf_1,\hat\acf_2,\hat\acf_3)\Tr$.
The factor $\acvs_0$ cancels because it multiplies \emph{both} sides equally;
hence $\hat{\boldsymbol\phi}=\bar\Gamma_3^{-1}\hat{\boldsymbol\acf}_3$ depends only
on the correlations. With $\hat\acf_0=1$ the explicit system is
\[
\begin{pmatrix}1 & 0.8 & 0.6\\ 0.8 & 1 & 0.8\\ 0.6 & 0.8 & 1\end{pmatrix}
\begin{pmatrix}\hat\phi_1\\ \hat\phi_2\\ \hat\phi_3\end{pmatrix}
=\begin{pmatrix}0.8\\ 0.6\\ 0.4\end{pmatrix}.
\]

\textbf{(b)} Solve by elimination. Label the rows $R_1,R_2,R_3$.
\[
\begin{aligned}
R_1:&\quad \hat\phi_1+0.8\hat\phi_2+0.6\hat\phi_3=0.8,\\
R_2:&\quad 0.8\hat\phi_1+\hat\phi_2+0.8\hat\phi_3=0.6,\\
R_3:&\quad 0.6\hat\phi_1+0.8\hat\phi_2+\hat\phi_3=0.4 .
\end{aligned}
\]
Eliminate $\hat\phi_1$. Compute $R_2-0.8R_1$ and $R_3-0.6R_1$:
\[
\begin{aligned}
R_2':&\quad (1-0.64)\hat\phi_2+(0.8-0.48)\hat\phi_3=0.6-0.64
\;\Rightarrow\; 0.36\hat\phi_2+0.32\hat\phi_3=-0.04,\\
R_3':&\quad (0.8-0.48)\hat\phi_2+(1-0.36)\hat\phi_3=0.4-0.48
\;\Rightarrow\; 0.32\hat\phi_2+0.64\hat\phi_3=-0.08 .
\end{aligned}
\]
From $R_3'$, divide by $0.32$: $\hat\phi_2+2\hat\phi_3=-0.25$. From $R_2'$,
divide by $0.04$: $9\hat\phi_2+8\hat\phi_3=-1$. Substitute
$\hat\phi_2=-0.25-2\hat\phi_3$:
\[
9(-0.25-2\hat\phi_3)+8\hat\phi_3=-1
\;\Rightarrow\;
-2.25-10\hat\phi_3=-1
\;\Rightarrow\;
\hat\phi_3=\frac{-1.25}{10}=-\tfrac18=-0.125 .
\]
Then $\hat\phi_2=-0.25-2(-0.125)=-0.25+0.25=0$, and from $R_1$,
$\hat\phi_1=0.8-0.8(0)-0.6(-0.125)=0.8+0.075=0.875=\tfrac78$. Thus
\[
\boxed{\;\hat\phi_1=\tfrac78=0.875,\qquad \hat\phi_2=0,\qquad \hat\phi_3=-\tfrac18=-0.125\;}
\]
\textit{Check.}
$R_1:\ 0.875+0.8(0)+0.6(-0.125)=0.875-0.075=0.8$.\ \checkmark\;
$R_2:\ 0.8(0.875)+0+0.8(-0.125)=0.7-0.1=0.6$.\ \checkmark\;
$R_3:\ 0.6(0.875)+0.8(0)+(-0.125)=0.525-0.125=0.4$.\ \checkmark

\textbf{(c)} Using $\hat\sigma^2=\hat\acvs_0\big(1-\sum_j\hat\phi_j\hat\acf_j\big)$
with $\hat\acvs_0=10$,
\[
1-\hat\phi_1\hat\acf_1-\hat\phi_2\hat\acf_2-\hat\phi_3\hat\acf_3
=1-\tfrac78(0.8)-0(0.6)-\big(-\tfrac18\big)(0.4)
=1-0.7+0.05=0.35=\tfrac{7}{20},
\]
hence $\hat\sigma^2=10\cdot\tfrac{7}{20}=\tfrac72=3.5$.
(The series length $n=150$ plays no role in the point estimates; it would only
enter standard errors.)

% ---------------------------------------------------------------
\solproblem{5}
\statementstart
This problem treats the two faces of least squares.

\textbf{(a)} \emph{(MA via residual reconstruction.)} A mean-zero $\mathrm{MA}(1)$ process $X_t=\eps_t-\theta\eps_{t-1}$ is observed at $X_1=0$, $X_2=2$, $X_3=-1$. Reconstruct the innovations $\hat\eps_1,\hat\eps_2,\hat\eps_3$ as functions of $\theta$ (state your initialisation), write the least-squares criterion $S(\theta)$, and find $\hat\theta$. Is the fit invertible? \hfill(2.5 pts)

\textbf{(b)} \emph{(Forward least squares for $\mathrm{AR}(2)$.)} For the five observations $x_1,\dots,x_5=1,\,2,\,0,\,-1,\,-2$ (already mean-zero), write the forward design $\mathbf X_F=F\boldsymbol\phi+\boldsymbol\eps$ for an $\mathrm{AR}(2)$ fit, form the normal equations $F\Tr F\,\hat{\boldsymbol\phi}_F=F\Tr\mathbf X_F$, and solve for $\hat{\boldsymbol\phi}_F$. Then report $\hat\sigma_F^2$ using the correct degrees of freedom. \hfill(2 pts)

\textbf{(c)} State, in one sentence each, (i) why the \emph{backward} least-squares estimator $\hat{\boldsymbol\phi}_B$ estimates the \emph{same} coefficients $\boldsymbol\phi$, and (ii) whether a least-squares AR fit is guaranteed to be stationary. \hfill(0.5 pt)
\statementend

\textbf{(a)} The innovations are latent, so reconstruct them from the model
equation $\eps_t=X_t+\theta\eps_{t-1}$. Initialise the pre-sample innovation to
zero, $\eps_0=0$ (take $t=0$ as the start). Then
\[
\hat\eps_1=X_1+\theta\eps_0=0,\qquad
\hat\eps_2=X_2+\theta\hat\eps_1=2,\qquad
\hat\eps_3=X_3+\theta\hat\eps_2=-1+2\theta .
\]
The least-squares criterion is
\[
S(\theta)=\sum_{t=1}^3\hat\eps_t^{\,2}=0^2+2^2+(2\theta-1)^2=4+(2\theta-1)^2 .
\]
Only the last term depends on $\theta$ (because $X_1=0$ froze
$\hat\eps_1=0,\hat\eps_2=2$), and it is minimised when $2\theta-1=0$. Hence
\[
\boxed{\;\hat\theta=\tfrac12\;}\qquad\text{with } S(\hat\theta)=4 .
\]
Since $\abs{\hat\theta}=\tfrac12<1$ (equivalently $\Theta(z)=1-\tfrac12 z$ has
root $z=2$ outside the unit circle), the fit is \textbf{invertible}.

\textbf{(b)} For $\mathrm{AR}(2)$ ($p=2$) the forward regression uses the
observable times $t=3,4,5$, regressing each $x_t$ on $(x_{t-1},x_{t-2})$:
\[
\mathbf X_F=\begin{pmatrix}x_3\\ x_4\\ x_5\end{pmatrix}=\begin{pmatrix}0\\ -1\\ -2\end{pmatrix},
\qquad
F=\begin{pmatrix}x_2 & x_1\\ x_3 & x_2\\ x_4 & x_3\end{pmatrix}
=\begin{pmatrix}2 & 1\\ 0 & 2\\ -1 & 0\end{pmatrix}.
\]
Form the normal equations $F\Tr F\,\hat{\boldsymbol\phi}_F=F\Tr\mathbf X_F$:
\[
F\Tr F=\begin{pmatrix}2 & 0 & -1\\ 1 & 2 & 0\end{pmatrix}
\begin{pmatrix}2 & 1\\ 0 & 2\\ -1 & 0\end{pmatrix}
=\begin{pmatrix}5 & 2\\ 2 & 5\end{pmatrix},
\qquad
F\Tr\mathbf X_F=\begin{pmatrix}2 & 0 & -1\\ 1 & 2 & 0\end{pmatrix}
\begin{pmatrix}0\\ -1\\ -2\end{pmatrix}
=\begin{pmatrix}2\\ -2\end{pmatrix}.
\]
With $\det(F\Tr F)=25-4=21$,
\[
\hat{\boldsymbol\phi}_F=\frac{1}{21}\begin{pmatrix}5 & -2\\ -2 & 5\end{pmatrix}
\begin{pmatrix}2\\ -2\end{pmatrix}
=\frac{1}{21}\begin{pmatrix}10+4\\ -4-10\end{pmatrix}
=\frac{1}{21}\begin{pmatrix}14\\ -14\end{pmatrix}
=\begin{pmatrix}\tfrac23\\[2pt] -\tfrac23\end{pmatrix}.
\]
Residuals $r_t=x_t-\hat\phi_1x_{t-1}-\hat\phi_2x_{t-2}$ for $t=3,4,5$:
\[
\begin{aligned}
r_3&=0-\tfrac23(2)+\tfrac23(1)=-\tfrac23,\\
r_4&=-1-\tfrac23(0)+\tfrac23(2)=-1+\tfrac43=\tfrac13,\\
r_5&=-2-\tfrac23(-1)+\tfrac23(0)=-2+\tfrac23=-\tfrac43,
\end{aligned}
\qquad
\norm{\mathbf X_F-F\hat{\boldsymbol\phi}_F}^2=\tfrac49+\tfrac19+\tfrac{16}{9}=\tfrac{21}{9}=\tfrac73 .
\]
The degrees of freedom are $n-2p=5-4=1$, so
\[
\hat\sigma_F^2=\frac{\norm{\mathbf X_F-F\hat{\boldsymbol\phi}_F}^2}{n-2p}
=\frac{7/3}{1}=\tfrac73\approx2.33 .
\]
So $\hat{\boldsymbol\phi}_F=(\tfrac23,-\tfrac23)\Tr$, $\hat\sigma_F^2=\tfrac73$.

\textbf{(c)} (i) A stationary Gaussian $\mathrm{AR}(p)$ run \emph{backwards} in
time is the \emph{same} $\mathrm{AR}(p)$ model (its ACVS is symmetric,
$\acvs_{-\tau}=\acvs_\tau$, so forward and reversed processes have identical
second-order structure), so regressing each point on its \emph{future} lags
estimates the very same coefficients $\boldsymbol\phi$. (ii) \textbf{No}: the
OLS objective only minimises the residual sum of squares and places no
constraint on the root moduli, so $\hat\Phi(z)$ may have roots inside the unit
circle. (Here it happens to be stationary --- $\Phi(z)=1-\tfrac23z+\tfrac23z^2$
has roots $z=\tfrac{1\pm\sqrt{-5}}{2}$ of modulus $\sqrt{3/2}\approx1.22>1$ ---
but this is not guaranteed in general; Yule--Walker, by contrast, always yields a
stationary fit.)

% ---------------------------------------------------------------
\solproblem{5}
\statementstart
\textbf{(Revision of Week 3.)} Consider the $\mathrm{ARMA}(2,1)$ process
\[
\Big(1-\tfrac34\Bsh+\tfrac18\Bsh^2\Big)X_t=\Big(1+\tfrac13\Bsh\Big)\eps_t .
\]

\textbf{(a)} Factor the AR polynomial $\Phi(z)=1-\tfrac34 z+\tfrac18 z^2$, locate its roots, and confirm the process is causal. Then check invertibility from $\Theta(z)=1+\tfrac13 z$. \hfill(2 pts)

\textbf{(b)} Using a \emph{partial-fraction} expansion of $G(z)=\Theta(z)/\Phi(z)$, find a closed form for the $\mathrm{MA}(\infty)$ weights $\psi_j$ in $X_t=\sum_{j\ge0}\psi_j\eps_{t-j}$, and report $\psi_0,\psi_1,\psi_2$. \hfill(3 pts)
\statementend

\textbf{(a)} The AR polynomial is
$\Phi(z)=1-\tfrac34 z+\tfrac18 z^2$. Multiplying by $8$, $z^2-6z+8=(z-2)(z-4)$,
so
\[
\Phi(z)=\tfrac18(z-2)(z-4)=\Big(1-\tfrac12 z\Big)\Big(1-\tfrac14 z\Big),
\]
with roots $z=2$ and $z=4$. Both satisfy $\abs{z}>1$, i.e.\ lie outside the unit
circle, so the process is \textbf{causal} (stationary). The MA polynomial
$\Theta(z)=1+\tfrac13 z$ has root $z=-3$, $\abs{-3}=3>1$ outside the unit circle,
so the process is also \textbf{invertible} (this guarantees $G=\Theta/\Phi$ is
well defined on $\abs{z}\le1$).

\textbf{(b)} Write $X_t=G(\Bsh)\eps_t$ with
\[
G(z)=\frac{1+\tfrac13 z}{\big(1-\tfrac12 z\big)\big(1-\tfrac14 z\big)} .
\]
Partial-fraction in the reciprocal-root form
\[
G(z)=\frac{A}{1-\tfrac12 z}+\frac{B}{1-\tfrac14 z},
\qquad
1+\tfrac13 z=A\Big(1-\tfrac14 z\Big)+B\Big(1-\tfrac12 z\Big).
\]
Evaluate at the poles $z=2$ and $z=4$:
\[
\begin{aligned}
z=2:\quad &1+\tfrac23=A\Big(1-\tfrac12\Big)=\tfrac12 A
&&\Rightarrow\quad A=2\cdot\tfrac53=\tfrac{10}{3},\\
z=4:\quad &1+\tfrac43=B\Big(1-2\Big)=-B
&&\Rightarrow\quad B=-\tfrac73 .
\end{aligned}
\]
Each term expands as a geometric series convergent on the unit disk (since
$\abs{z}\le1<2,4$):
\[
\frac{A}{1-\tfrac12 z}=A\sum_{j\ge0}\Big(\tfrac12\Big)^j z^j,
\qquad
\frac{B}{1-\tfrac14 z}=B\sum_{j\ge0}\Big(\tfrac14\Big)^j z^j .
\]
Reading off the coefficient of $z^j$,
\[
\boxed{\;\psi_j=\frac{10}{3}\Big(\tfrac12\Big)^{j}-\frac{7}{3}\Big(\tfrac14\Big)^{j},\qquad j\ge0\;}
\]
which already gives $\psi_0=\tfrac{10}{3}-\tfrac73=1$. In particular
\[
\begin{aligned}
\psi_0&=\tfrac{10}{3}-\tfrac73=1,\\
\psi_1&=\tfrac{10}{3}\cdot\tfrac12-\tfrac73\cdot\tfrac14=\tfrac53-\tfrac{7}{12}=\tfrac{20-7}{12}=\tfrac{13}{12},\\
\psi_2&=\tfrac{10}{3}\cdot\tfrac14-\tfrac73\cdot\tfrac1{16}=\tfrac56-\tfrac{7}{48}=\tfrac{40-7}{48}=\tfrac{33}{48}=\tfrac{11}{16}.
\end{aligned}
\]
\textit{Check (coefficient matching).} Match powers of $z$ in
$\big(1-\tfrac34 z+\tfrac18 z^2\big)\sum_j\psi_j z^j=1+\tfrac13 z$:
\[
\psi_0=1,\qquad
\psi_1-\tfrac34\psi_0=\tfrac13\ \Rightarrow\ \psi_1=\tfrac34+\tfrac13=\tfrac{13}{12},
\]
and $\psi_j=\tfrac34\psi_{j-1}-\tfrac18\psi_{j-2}$ for $j\ge2$, giving
$\psi_2=\tfrac34\cdot\tfrac{13}{12}-\tfrac18=\tfrac{13}{16}-\tfrac18=\tfrac{11}{16}$,
matching exactly.

\end{document}
